% `tex.web` §435 and §1236, two numbers this engine took quietly.
%
% §435 `scan_four_bit_int` — a math family is 0..15 and nothing else:
%
%     procedure scan_four_bit_int;
%     begin scan_int;
%     if (cur_val<0)or(cur_val>15) then
%       begin print_err("Bad number");
%       help2("Since I expected to read a number between 0 and 15,")
%         ("I changed this one to zero."); int_error(cur_val); cur_val:=0;
%       end;
%     end;
%
% We took the number modulo sixteen, so `\textfont16` quietly meant
% `\textfont0` and `\scriptfont-1` meant `\scriptfont15`.  Both are errors,
% and both read as nought.  §91's `int_error` puts the offending value on the
% top line.
%
% §1236, on `\multiply` and `\divide`:
%
%     arith_error:=false;
%     …
%     if arith_error then
%       begin print_err("Arithmetic overflow");
%       help2("I can't carry out that multiplication or division,")
%       ("since the result is out of range.");
%       … error; return;
%       end
%
% The `return` is the point: the *old value stays*.  TeX declines the
% assignment rather than storing a wrapped number, which is the same shape as
% §1243's `\spacefactor`.  We wrapped and stored, so
% `\count0=100000 \multiply\count0 by 100000` left a negative number in the
% register and said nothing at all.  An integer is bounded at `07777777777`
% and a dimension at `max_dimen`.
\catcode`\{=1 \catcode`\}=2
\font\r=cmr10 \r \hsize100pt \parindent0pt
\textfont16=\r
\scriptfont-1=\r
\showthe\textfont16
\dimen0=16383.99998pt \multiply\dimen0 by 3 \showthe\dimen0
\count0=100000 \multiply\count0 by 100000 \showthe\count0
\count1=6 \multiply\count1 by 7 \showthe\count1
\skip0=1pt plus 2pt \multiply\skip0 by 3 \showthe\skip0
\end
